\documentclass[11pt]{article}
\usepackage[paperwidth=210mm,paperheight=297mm,textwidth=182mm,top=20mm,bottom=22mm]{geometry}
\usepackage{amsmath,amssymb,newtxmath}
\usepackage[no-math]{fontspec}
\setmainfont{Times New Roman}
\usepackage{xeCJK}
\setCJKmainfont[AutoFakeBold=2]{SimSun}
\usepackage{graphicx,booktabs,array,hyperref,placeins}
\hypersetup{hidelinks,unicode=true}
\raggedbottom
\setlength{\emergencystretch}{2em}
\renewcommand{\arraystretch}{1.15}
\usepackage{seqsplit}
\begin{document}
\title{Litz 线分层建模与数值验证：完整研究报告 · 12 结论}
\author{缪雨轩\\淮南师范学院}
\date{}
\maketitle
\noindent 合并稿对应章节 / Chapter in merged manuscript: \href{report_zh-CN_single.pdf#chapter-12}{12}\par
\FloatBarrier
\section*{12 结论}
\hypertarget{chapter-12}{}
\pdfbookmark[0]{12 结论}{bookmark-chapter-12}
第一，分层验证框架把共同物理内容与各层缺口分别记录，使解析、二维、纵向及三维外场结果能够在可比较范围内使用。第二，共端增广系统给出了指定线性模型内的复电流均流判据，避免用预设均流替代验证。第三，有限长度纵向 PEEC 的已运行长度与积分敏感性达到专项门槛，其末档电阻与本档等面积解析参照相符；该结论继续受纵向基函数范围限制。


640 股二维场、固定外场长方体 RT0 和 CPU/GPU 同模型对照形成可重绘的冻结证据。实际换位束的三维电磁排名、热模型和实验仍未收口；失败路线保留触发条件、原始结果和后续所需能力。


\end{document}
